Input to the Well (WEL) Package is read from the file that has type ``WEL6'' in the Name File.  Any number of WEL Packages can be specified for a single groundwater flow model.

Well input can be specified using lists or grid arrays.  Grid array-based input for the WEL package is activated by providing READARRAYGRID within the OPTIONS block.  Instructions for specifying list-based well input are provided in the previous section.

When grid array-based input is used for the well package, the DIMENSIONS block is optional.  The input array size is determined from the model grid.  Cells that define no data should contain the DNODATA (3.0E+30) value.  When defined, the MAXBOUND dimension represents the maximum number of defined well (non-DNODATA) cells found in any stress period.  Defining MAXBOUND can reduce the memory footprint of the grid array-based WEL package, especially when the input grid arrays are sparse and the model grid is large.

\vspace{5mm}
\subsubsection{Structure of Blocks}
\vspace{5mm}

\noindent \textit{FOR EACH SIMULATION}
\lstinputlisting[style=blockdefinition]{./mf6ivar/tex/gwf-welg-options.dat}
\lstinputlisting[style=blockdefinition]{./mf6ivar/tex/gwf-welg-dimensions.dat}
\vspace{5mm}
\noindent \textit{FOR ANY STRESS PERIOD}
\lstinputlisting[style=blockdefinition]{./mf6ivar/tex/gwf-welg-period.dat}
\packageperioddescriptiongridarray{well}

\vspace{5mm}
\subsubsection{Explanation of Variables}
\begin{description}
\input{./mf6ivar/tex/gwf-welg-desc.tex}
\end{description}

\vspace{5mm}
\subsubsection{Example Input File}
\lstinputlisting[style=inputfile]{./mf6ivar/examples/gwf-welg-example.dat}

\vspace{5mm}
\subsubsection{Available observation types}
Well Package observations include the simulated well rates (\texttt{wel}), the well discharge that is available for the MVR package (\texttt{to-mvr}), and the reduction in the specified \texttt{q} when the \texttt{AUTO\_FLOW\_REDUCE} option is enabled. The data required for each WEL Package observation type is defined in table~\ref{table:gwf-welgobstype}. The sum of \texttt{wel} and \texttt{to-mvr} is equal to the simulated well discharge rate, which may be less than the specified \texttt{q} if the \texttt{AUTO\_FLOW\_REDUCE} option is enabled. The \texttt{DNODATA} value is returned if the \texttt{wel-reduction} observation is specified but the \texttt{AUTO\_FLOW\_REDUCE} option is not enabled. Negative and positive values for an observation represent a loss from and gain to the GWF model, respectively.

\begin{longtable}{p{2cm} p{2.75cm} p{2cm} p{1.25cm} p{7cm}}
\caption{Available WEL Package observation types} \tabularnewline

\hline
\hline
\textbf{Stress Package} & \textbf{Observation type} & \textbf{ID} & \textbf{ID2} & \textbf{Description} \\
\hline
\endhead

\hline
\endfoot

\input{../Common/gwf-welobs.tex}
\label{table:gwf-welgobstype}
\end{longtable}

\vspace{5mm}
\subsubsection{Example Observation Input File}
\lstinputlisting[style=inputfile]{./mf6ivar/examples/gwf-welg-example-obs.dat}
